% GENERATED FILE. Edit canonical lessons, then run npm run generate:books.
% canonical-source-hash: fnv1a64:73232d284d6ce824
% canonical-lessons: KA-C388-naitika, KA-C388-tyaga, KA-C388-manassaksi, KA-C388-hita, KA-C388-khandisu

\chapter{Right and Wrong}
\label{ch:ka-right-and-wrong}
\hlchaptermodality{388}

\begin{chapteropening}
\textbf{By the end of this chapter:} \emph{I can say ethical, a sacrifice, a conscience, well-being and to condemn.}

Complete the last lesson of chapter 388: I can say ethical, a sacrifice, a conscience, well-being and to condemn.
\end{chapteropening}

\section[\texorpdfstring{\kn{ನೈತಿಕ} (naitika)}{ನೈತಿಕ (naitika)}]{\kn{ನೈತಿಕ} (naitika) — ethical; moral}

\label{lesson:KA-C388-naitika}



\begin{quote}
Before the new one: say the Kannada for a priority, then the Kannada for a value.
\end{quote}

\subsection*{You'll want to know: \kn{ನೈತಿಕ}}
\textbf{\kn{ನೈತಿಕ}} — \emph{naitika} — \textquotedblleft{}ethical; moral\textquotedblright{}.

\begin{quote}
\kn{ಇದು} \kn{ಹಣದ} \kn{ಪ್ರಶ್ನೆ} \kn{ಅಲ್ಲ}, \kn{ನೈತಿಕ} \kn{ಪ್ರಶ್ನೆ}.
\end{quote}

\begin{quote}
\emph{idu haṇada praśne alla, naitika praśne.}
\end{quote}

\begin{quote}
This is not a question of money but a moral one.
\end{quote}

\kn{ನೈತಿಕತೆ}, morality.

\begin{cousinweb}
\textbf{\kn{ನೈತಿಕ}} (\emph{naitika}) is built on Sanskrit \emph{nīti}, \textquotedblleft{}conduct, right behaviour\textquotedblright{} (CDIAL 7546), with the adjective ending -\emph{ika}.
\end{cousinweb}

\begin{grammarlens}[title={What you've built}]
One more word for giving reasons and opinions.
\end{grammarlens}

\subsection*{Your turn}
\begin{itemize}
\raggedright
  \item \emph{Say it:} \emph{naitika}
  \item \emph{Say it:} \emph{naitika}, once more
  \item \emph{Say it:} \emph{idu haṇada praśne alla, naitika praśne.} — This is not a question of money but a moral one.
  \item \emph{Recall:} say the Kannada for a priority, then the Kannada for a value, then say \emph{naitika} again
\end{itemize}

\begin{tcolorbox}[breakable,colback=teal!4,colframe=teal!35!black,title={Before you move on}]
What does \kn{ನೈತಿಕ} mean? (\textquotedblleft{}Ethical; moral\textquotedblright{}.) Say it once more.
\end{tcolorbox}
\section[\texorpdfstring{\kn{ತ್ಯಾಗ} (tyāga)}{ತ್ಯಾಗ (tyāga)}]{\kn{ತ್ಯಾಗ} (tyāga) — a sacrifice}

\label{lesson:KA-C388-tyaga}



\begin{quote}
Before the new one: say the Kannada for a value, then the Kannada for ethical.
\end{quote}

\subsection*{You'll want to know: \kn{ತ್ಯಾಗ}}
\textbf{\kn{ತ್ಯಾಗ}} — \emph{tyāga} — \textquotedblleft{}a sacrifice\textquotedblright{}.

\begin{quote}
\kn{ತಾಯಿ} \kn{ಮಕ್ಕಳಿಗಾಗಿ} \kn{ತುಂಬಾ} \kn{ತ್ಯಾಗ} \kn{ಮಾಡಿದಳು}.
\end{quote}

\begin{quote}
\emph{tāyi makkaḷigāgi tumbā tyāga māḍidaḷu.}
\end{quote}

\begin{quote}
The mother made many sacrifices for her children.
\end{quote}

\kn{ತ್ಯಾಗ} \kn{ಮಾಡು}, to make a sacrifice.

\begin{cousinweb}
\textbf{\kn{ತ್ಯಾಗ}} (\emph{tyāga}) is borrowed from Sanskrit \emph{tyāga}, \textquotedblleft{}leaving, giving up\textquotedblright{} (CDIAL 5985).
\end{cousinweb}

\begin{grammarlens}[title={What you've built}]
One more word for giving reasons and opinions.
\end{grammarlens}

\subsection*{Your turn}
\begin{itemize}
\raggedright
  \item \emph{Say it:} \emph{tyāga}
  \item \emph{Say it:} \emph{tyāga}, once more
  \item \emph{Say it:} \emph{tāyi makkaḷigāgi tumbā tyāga māḍidaḷu.} — The mother made many sacrifices for her children.
  \item \emph{Recall:} say the Kannada for a value, then the Kannada for ethical, then say \emph{tyāga} again
\end{itemize}

\begin{tcolorbox}[breakable,colback=teal!4,colframe=teal!35!black,title={Before you move on}]
What does \kn{ತ್ಯಾಗ} mean? (\textquotedblleft{}A sacrifice\textquotedblright{}.) Say it once more.
\end{tcolorbox}
\section[\texorpdfstring{\kn{ಮನಸ್ಸಾಕ್ಷಿ} (manassākṣi)}{ಮನಸ್ಸಾಕ್ಷಿ (manassākṣi)}]{\kn{ಮನಸ್ಸಾಕ್ಷಿ} (manassākṣi) — a conscience}

\label{lesson:KA-C388-manassaksi}



\begin{quote}
Before the new one: say the Kannada for ethical, then the Kannada for a sacrifice.
\end{quote}

\subsection*{You'll want to know: \kn{ಮನಸ್ಸಾಕ್ಷಿ}}
\textbf{\kn{ಮನಸ್ಸಾಕ್ಷಿ}} — \emph{manassākṣi} — \textquotedblleft{}a conscience\textquotedblright{}.

\begin{quote}
\kn{ನನ್ನ} \kn{ಮನಸ್ಸಾಕ್ಷಿ} \kn{ಸುಳ್ಳು} \kn{ಹೇಳಲು} \kn{ಬಿಡಲಿಲ್ಲ}.
\end{quote}

\begin{quote}
\emph{nanna manassākṣi suḷḷu hēḷalu biḍalilla.}
\end{quote}

\begin{quote}
My conscience did not let me lie.
\end{quote}

\kn{ಮನಸ್ಸು} (mind) + \kn{ಸಾಕ್ಷಿ} (witness).

\begin{cousinweb}
\textbf{\kn{ಮನಸ್ಸಾಕ್ಷಿ}} (\emph{manassākṣi}) is built from Sanskrit \emph{manas}, \textquotedblleft{}mind\textquotedblright{} (CDIAL 9822), and \emph{sākṣin}, \textquotedblleft{}a witness\textquotedblright{} (CDIAL 13321): the mind as witness.
\end{cousinweb}

\begin{grammarlens}[title={What you've built}]
One more word for giving reasons and opinions.
\end{grammarlens}

\subsection*{Your turn}
\begin{itemize}
\raggedright
  \item \emph{Say it:} \emph{manassākṣi}
  \item \emph{Say it:} \emph{manassākṣi}, once more
  \item \emph{Say it:} \emph{nanna manassākṣi suḷḷu hēḷalu biḍalilla.} — My conscience did not let me lie.
  \item \emph{Recall:} say the Kannada for ethical, then the Kannada for a sacrifice, then say \emph{manassākṣi} again
\end{itemize}

\begin{tcolorbox}[breakable,colback=teal!4,colframe=teal!35!black,title={Before you move on}]
What does \kn{ಮನಸ್ಸಾಕ್ಷಿ} mean? (\textquotedblleft{}A conscience\textquotedblright{}.) Say it once more.
\end{tcolorbox}
\section[\texorpdfstring{\kn{ಹಿತ} (hita)}{ಹಿತ (hita)}]{\kn{ಹಿತ} (hita) — well-being; good}

\label{lesson:KA-C388-hita}



\begin{quote}
Before the new one: say the Kannada for a sacrifice, then the Kannada for a conscience.
\end{quote}

\subsection*{You'll want to know: \kn{ಹಿತ}}
\textbf{\kn{ಹಿತ}} — \emph{hita} — \textquotedblleft{}well-being; good\textquotedblright{}.

\begin{quote}
\kn{ಊರಿನ} \kn{ಹಿತಕ್ಕಾಗಿ} \kn{ಎಲ್ಲರೂ} \kn{ಕೆಲಸ} \kn{ಮಾಡಿದರು}.
\end{quote}

\begin{quote}
\emph{ūrina hitakkāgi ellarū kelasa māḍidaru.}
\end{quote}

\begin{quote}
Everyone worked for the good of the village.
\end{quote}

\kn{ಹಿತವಾದ}, pleasant; \kn{ಅಹಿತ}, harmful.

\begin{cousinweb}
\textbf{\kn{ಹಿತ}} (\emph{hita}) is borrowed from Sanskrit \emph{hita}, \textquotedblleft{}beneficial, good, welfare\textquotedblright{} (CDIAL 14091).
\end{cousinweb}

\begin{grammarlens}[title={What you've built}]
One more word for giving reasons and opinions.
\end{grammarlens}

\subsection*{Your turn}
\begin{itemize}
\raggedright
  \item \emph{Say it:} \emph{hita}
  \item \emph{Say it:} \emph{hita}, once more
  \item \emph{Say it:} \emph{ūrina hitakkāgi ellarū kelasa māḍidaru.} — Everyone worked for the good of the village.
  \item \emph{Recall:} say the Kannada for a sacrifice, then the Kannada for a conscience, then say \emph{hita} again
\end{itemize}

\begin{tcolorbox}[breakable,colback=teal!4,colframe=teal!35!black,title={Before you move on}]
What does \kn{ಹಿತ} mean? (\textquotedblleft{}Well-being; good\textquotedblright{}.) Say it once more.
\end{tcolorbox}
\section[\texorpdfstring{\kn{ಖಂಡಿಸು} (khaṇḍisu)}{ಖಂಡಿಸು (khaṇḍisu)}]{\kn{ಖಂಡಿಸು} (khaṇḍisu) — to condemn}

\label{lesson:KA-C388-khandisu}



\begin{quote}
Before the new one: say the Kannada for a conscience, then the Kannada for well-being.
\end{quote}

\subsection*{You'll want to know: \kn{ಖಂಡಿಸು}}
\textbf{\kn{ಖಂಡಿಸು}} — \emph{khaṇḍisu} — \textquotedblleft{}to condemn\textquotedblright{}.

\begin{quote}
\kn{ಎಲ್ಲರೂ} \kn{ಆ} \kn{ಕೆಟ್ಟ} \kn{ಕೆಲಸವನ್ನು} \kn{ಖಂಡಿಸಿದರು}.
\end{quote}

\begin{quote}
\emph{ellarū ā keṭṭa kelasavannu khaṇḍisidaru.}
\end{quote}

\begin{quote}
Everyone condemned that wicked act.
\end{quote}

Past (I …): \kn{ಖಂಡಿಸಿದೆ}, \emph{khaṇḍiside}.

What you condemn takes -\kn{ಅನ್ನು}. \kn{ಖಂಡನೆ}, condemnation.

\begin{cousinweb}
\textbf{\kn{ಖಂಡಿಸು}} (\emph{khaṇḍisu}) is built on Sanskrit \emph{khaṇḍana}, \textquotedblleft{}breaking, refuting\textquotedblright{} (CDIAL 3796), with the Kannada verb ending -\emph{isu}.
\end{cousinweb}

\begin{grammarlens}[title={What you've built}]
One more word for giving reasons and opinions.
\end{grammarlens}

\subsection*{Your turn}
\begin{itemize}
\raggedright
  \item \emph{Say it:} \emph{khaṇḍisu}
  \item \emph{Say it:} \emph{khaṇḍisu}, once more
  \item \emph{Say it:} \emph{ellarū ā keṭṭa kelasavannu khaṇḍisidaru.} — Everyone condemned that wicked act.
  \item \emph{Recall:} say the Kannada for a conscience, then the Kannada for well-being, then say \emph{khaṇḍisu} again
\end{itemize}

\begin{tcolorbox}[breakable,colback=teal!4,colframe=teal!35!black,title={Before you move on}]
What does \kn{ಖಂಡಿಸು} mean? (\textquotedblleft{}To condemn\textquotedblright{}.) Say it once more.
\end{tcolorbox}
